%!TEX root = ../main.tex
\chapter{CI/CD pour le Machine Learning}
\label{ch:cicd}

\section{Introduction}

Dans le d\'eveloppement logiciel traditionnel, les pratiques de CI/CD (Int\'egration Continue / D\'eploiement Continu) ont r\'evolutionn\'e la mani\`ere dont les \'equipes livrent du code. En Machine Learning, ces pratiques deviennent encore plus critiques~: il ne s'agit plus seulement de tester du code, mais \'egalement de valider des donn\'ees, des mod\`eles et des pipelines d'entra\^inement complets.

\begin{definition}[Int\'egration Continue (CI)]
L'int\'egration continue est une pratique de d\'eveloppement o\`u les d\'eveloppeurs fusionnent fr\'equemment leurs modifications dans une branche principale. Chaque fusion d\'eclenche une compilation et des tests automatis\'es pour d\'etecter les erreurs le plus t\^ot possible.
\end{definition}

\begin{definition}[Livraison Continue (CD)]
La livraison continue \'etend la CI en automatisant le processus de mise en production. On distingue~:
\begin{itemize}
    \item \textbf{Continuous Delivery}~: le code est pr\^et \`a \^etre d\'eploy\'e \`a tout moment, mais le d\'eploiement reste manuel.
    \item \textbf{Continuous Deployment}~: chaque modification valid\'ee est automatiquement d\'eploy\'ee en production.
\end{itemize}
\end{definition}

%% ----------------------------------------------------------------
\section{CI/CD pour le ML~: qu'est-ce qui change~?}
\label{sec:cicd-ml-specifics}

Le ML introduit des d\'efis suppl\'ementaires par rapport au d\'eveloppement logiciel classique.

\begin{attention}[title=\textbf{Sp\'ecificit\'es du CI/CD en ML}]
Contrairement au logiciel traditionnel, un pipeline ML doit valider~:
\begin{enumerate}
    \item \textbf{Le code} -- qualit\'e, linting, typage
    \item \textbf{Les donn\'ees} -- sch\'ema, distribution, int\'egrit\'e
    \item \textbf{Le mod\`ele} -- performance, \'equit\'e, robustesse
    \item \textbf{L'infrastructure} -- reproductibilit\'e, ressources
\end{enumerate}
Un changement de donn\'ees peut casser un mod\`ele \emph{sans qu'aucune ligne de code n'ait chang\'e}.
\end{attention}

\begin{formules}[title=\textbf{Diff\'erences CI/CD classique vs ML}]
\begin{center}
\begin{tabular}{lll}
\toprule
\textbf{Aspect} & \textbf{CI/CD classique} & \textbf{CI/CD ML} \\
\midrule
Entr\'ee & Code source & Code + Donn\'ees + Mod\`ele \\
Tests & Unitaires, int\'egration & + validation donn\'ees, perf. mod\`ele \\
Artefact & Binaire, image Docker & + mod\`ele s\'erialis\'e, m\'etriques \\
D\'eclencheur & Commit, PR & + nouveau dataset, drift d\'etect\'e \\
Validation & Tests passent & + seuils de performance atteints \\
D\'eploiement & Blue/Green, Canary & + A/B testing, Shadow mode \\
\bottomrule
\end{tabular}
\end{center}
\end{formules}

\subsection{Strat\'egies de d\'eploiement ML}

\begin{definition}[A/B Testing pour les mod\`eles]
L'A/B testing en ML consiste \`a exposer simultan\'ement deux versions d'un mod\`ele \`a des sous-ensembles d'utilisateurs, puis \`a comparer les m\'etriques m\'etier (taux de clic, revenu, satisfaction) pour d\'eterminer quelle version est sup\'erieure. Cette approche n\'ecessite une infrastructure capable de router le trafic et de collecter des m\'etriques en temps r\'eel.
\end{definition}

\begin{definition}[Shadow Deployment]
Le \emph{shadow deployment} (ou \emph{dark launch}) consiste \`a d\'eployer un nouveau mod\`ele en parall\`ele du mod\`ele en production. Le nouveau mod\`ele re\c{c}oit le m\^eme trafic, mais ses pr\'edictions ne sont pas expos\'ees aux utilisateurs. Cela permet de comparer les performances sans risque.
\end{definition}

%% ----------------------------------------------------------------
\section{GitHub Actions~: workflow complet pour un projet ML}
\label{sec:github-actions}

GitHub Actions est l'outil de CI/CD int\'egr\'e \`a GitHub. Il utilise des fichiers YAML pour d\'efinir des workflows d\'eclench\'es par des \'ev\'enements (push, pull request, etc.).

\begin{codeblock}[title=\textbf{Workflow CI/CD complet pour un projet ML}]
\begin{minted}{yaml}
# .github/workflows/ml-pipeline.yml
name: ML Pipeline CI/CD

on:
  push:
    branches: [main, develop]
  pull_request:
    branches: [main]

env:
  PYTHON_VERSION: "3.11"
  DVC_REMOTE: s3://my-bucket/dvc-store

jobs:
  lint-and-format:
    runs-on: ubuntu-latest
    steps:
      - uses: actions/checkout@v4
      - uses: actions/setup-python@v5
        with:
          python-version: ${{ env.PYTHON_VERSION }}
      - name: Install linting tools
        run: pip install ruff black mypy
      - name: Check formatting (Black)
        run: black --check src/ tests/
      - name: Lint (Ruff)
        run: ruff check src/ tests/
      - name: Type checking (mypy)
        run: mypy src/ --ignore-missing-imports

  test:
    runs-on: ubuntu-latest
    needs: lint-and-format
    steps:
      - uses: actions/checkout@v4
      - uses: actions/setup-python@v5
        with:
          python-version: ${{ env.PYTHON_VERSION }}
      - name: Install dependencies
        run: |
          pip install -r requirements.txt
          pip install pytest pytest-cov
      - name: Run unit tests
        run: pytest tests/unit/ -v --cov=src/ --cov-report=xml
      - name: Run integration tests
        run: pytest tests/integration/ -v
      - name: Upload coverage
        uses: codecov/codecov-action@v3

  train-and-validate:
    runs-on: ubuntu-latest
    needs: test
    steps:
      - uses: actions/checkout@v4
      - uses: actions/setup-python@v5
        with:
          python-version: ${{ env.PYTHON_VERSION }}
      - name: Install DVC
        run: pip install dvc[s3]
      - name: Pull data with DVC
        run: dvc pull
        env:
          AWS_ACCESS_KEY_ID: ${{ secrets.AWS_ACCESS_KEY_ID }}
          AWS_SECRET_ACCESS_KEY: ${{ secrets.AWS_SECRET_KEY }}
      - name: Reproduce pipeline
        run: dvc repro
      - name: Validate model performance
        run: python scripts/validate_model.py
          --min-accuracy 0.85
          --max-latency-ms 100
          --check-drift

  deploy:
    runs-on: ubuntu-latest
    needs: train-and-validate
    if: github.ref == 'refs/heads/main'
    steps:
      - uses: actions/checkout@v4
      - name: Build Docker image
        run: docker build -t ml-model:${{ github.sha }} .
      - name: Push to registry
        run: |
          docker tag ml-model:${{ github.sha }} \
            ghcr.io/${{ github.repository }}/model:latest
          docker push ghcr.io/${{ github.repository }}/model:latest
      - name: Deploy to staging
        run: |
          kubectl set image deployment/ml-model \
            model=ghcr.io/${{ github.repository }}/model:latest
\end{minted}
\end{codeblock}

%% ----------------------------------------------------------------
\section{Pre-commit hooks pour la qualit\'e du code}
\label{sec:pre-commit}

Les \emph{pre-commit hooks} permettent d'ex\'ecuter des v\'erifications automatiques avant chaque commit, garantissant que le code qui entre dans le d\'ep\^ot respecte les standards de qualit\'e.

\begin{codeblock}[title=\textbf{Configuration \file{.pre-commit-config.yaml}}]
\begin{minted}{yaml}
repos:
  - repo: https://github.com/psf/black
    rev: 24.4.2
    hooks:
      - id: black
        language_version: python3.11

  - repo: https://github.com/astral-sh/ruff-pre-commit
    rev: v0.4.4
    hooks:
      - id: ruff
        args: [--fix, --exit-non-zero-on-fix]

  - repo: https://github.com/pre-commit/mirrors-mypy
    rev: v1.10.0
    hooks:
      - id: mypy
        additional_dependencies: [numpy, pandas-stubs]

  - repo: https://github.com/pre-commit/pre-commit-hooks
    rev: v4.6.0
    hooks:
      - id: check-yaml
      - id: check-json
      - id: check-added-large-files
        args: [--maxkb=500]
      - id: trailing-whitespace
      - id: end-of-file-fixer
\end{minted}
\end{codeblock}

\begin{bonnepratique}[title=\textbf{Installation et utilisation}]
\begin{minted}{bash}
# Installation
pip install pre-commit

# Initialiser les hooks dans le depot
pre-commit install

# Executer sur tous les fichiers (premiere fois)
pre-commit run --all-files
\end{minted}
Les hooks s'ex\'ecutent automatiquement \`a chaque \cli{git commit}. Si un hook \'echoue, le commit est bloqu\'e jusqu'\`a correction.
\end{bonnepratique}

\begin{remarque}
L'option \cli{--maxkb=500} dans \texttt{check-added-large-files} emp\^eche l'ajout accidentel de fichiers volumineux (datasets, mod\`eles s\'erialis\'es) dans le d\'ep\^ot Git. Utilisez DVC pour versionner ces fichiers \`a la place.
\end{remarque}

%% ----------------------------------------------------------------
\section{Validation de mod\`ele dans la CI}
\label{sec:model-validation-ci}

La validation automatis\'ee du mod\`ele est l'\'etape la plus sp\'ecifique au ML dans un pipeline CI/CD.

\begin{codeblock}[title=\textbf{Script de validation de mod\`ele}]
\begin{minted}{python}
# scripts/validate_model.py
import argparse
import json
import sys
from pathlib import Path

import joblib
import numpy as np
from scipy import stats
from sklearn.metrics import accuracy_score, f1_score


def load_metrics(path: str = "metrics.json") -> dict:
    """Charge les metriques du modele entraine."""
    with open(path) as f:
        return json.load(f)


def check_performance(metrics: dict, min_accuracy: float) -> bool:
    """Verifie que les metriques depassent les seuils."""
    accuracy = metrics.get("accuracy", 0.0)
    if accuracy < min_accuracy:
        print(f"ECHEC: accuracy={accuracy:.4f} < {min_accuracy}")
        return False
    print(f"OK: accuracy={accuracy:.4f} >= {min_accuracy}")
    return True


def check_data_drift(
    reference_path: str, current_path: str, threshold: float = 0.05
) -> bool:
    """Detecte le drift via le test de Kolmogorov-Smirnov."""
    ref = np.load(reference_path)
    cur = np.load(current_path)

    drift_detected = False
    for i in range(ref.shape[1]):
        stat, p_value = stats.ks_2samp(ref[:, i], cur[:, i])
        if p_value < threshold:
            print(f"DRIFT feature {i}: KS={stat:.4f}, p={p_value:.6f}")
            drift_detected = True

    return not drift_detected


def check_latency(model_path: str, sample_path: str,
                  max_latency_ms: float) -> bool:
    """Verifie la latence d'inference."""
    import time
    model = joblib.load(model_path)
    sample = np.load(sample_path)[:100]

    start = time.perf_counter()
    for x in sample:
        model.predict(x.reshape(1, -1))
    elapsed_ms = (time.perf_counter() - start) / len(sample) * 1000

    if elapsed_ms > max_latency_ms:
        print(f"ECHEC: latence={elapsed_ms:.2f}ms > {max_latency_ms}ms")
        return False
    print(f"OK: latence={elapsed_ms:.2f}ms <= {max_latency_ms}ms")
    return True


if __name__ == "__main__":
    parser = argparse.ArgumentParser()
    parser.add_argument("--min-accuracy", type=float, default=0.85)
    parser.add_argument("--max-latency-ms", type=float, default=100)
    parser.add_argument("--check-drift", action="store_true")
    args = parser.parse_args()

    metrics = load_metrics()
    checks = [check_performance(metrics, args.min_accuracy)]

    if args.check_drift:
        checks.append(check_data_drift(
            "data/reference.npy", "data/current.npy"
        ))

    checks.append(check_latency(
        "models/model.pkl", "data/current.npy", args.max_latency_ms
    ))

    if not all(checks):
        print("\nValidation ECHOUEE -- deploiement bloque.")
        sys.exit(1)
    print("\nValidation REUSSIE -- pret pour le deploiement.")
\end{minted}
\end{codeblock}

%% ----------------------------------------------------------------
\section{DVC dans la CI~: reproduire les pipelines automatiquement}
\label{sec:dvc-ci}

DVC (Data Version Control) permet de versionner les donn\'ees et de d\'efinir des pipelines reproductibles. Dans un contexte CI, DVC garantit que le pipeline est ex\'ecut\'e de mani\`ere d\'eterministe.

\begin{codeblock}[title=\textbf{Pipeline DVC \file{dvc.yaml}}]
\begin{minted}{yaml}
stages:
  preprocess:
    cmd: python src/preprocess.py
    deps:
      - src/preprocess.py
      - data/raw/
    outs:
      - data/processed/
    params:
      - preprocess.test_size
      - preprocess.seed

  train:
    cmd: python src/train.py
    deps:
      - src/train.py
      - data/processed/
    outs:
      - models/model.pkl
    params:
      - train.n_estimators
      - train.learning_rate
    metrics:
      - metrics.json:
          cache: false

  evaluate:
    cmd: python src/evaluate.py
    deps:
      - src/evaluate.py
      - models/model.pkl
      - data/processed/test.csv
    metrics:
      - eval_metrics.json:
          cache: false
    plots:
      - plots/confusion_matrix.csv:
          x: predicted
          y: actual
\end{minted}
\end{codeblock}

\begin{bonnepratique}[title=\textbf{Commandes DVC dans la CI}]
\begin{minted}{bash}
# Reproduire le pipeline (seules les etapes modifiees)
dvc repro

# Comparer les metriques avec la branche principale
dvc metrics diff main

# Afficher les parametres
dvc params diff main
\end{minted}
La commande \cli{dvc repro} n'ex\'ecute que les \'etapes dont les d\'ependances ont chang\'e, ce qui acc\'el\`ere consid\'erablement les pipelines CI.
\end{bonnepratique}

%% ----------------------------------------------------------------
\section{Architecture d'un pipeline CI/CD pour le ML}
\label{sec:cicd-architecture}

\begin{figure}[htbp]
\centering
\begin{tikzpicture}[
    node distance=1.8cm and 2.5cm,
    block/.style={rectangle, draw, rounded corners, fill=blue!10,
        minimum width=2.8cm, minimum height=1cm, align=center,
        font=\small},
    decision/.style={diamond, draw, fill=orange!15, aspect=2,
        inner sep=1pt, font=\small, align=center},
    arrow/.style={-{Stealth[length=2.5mm]}, thick},
    label/.style={font=\scriptsize, text=gray}
]

% Nodes
\node[block] (commit) {Commit /\\Pull Request};
\node[block, right=of commit] (lint) {Lint \&\\Formatage};
\node[block, right=of lint] (test) {Tests\\unitaires};
\node[block, below=of test] (data) {Validation\\donn\'ees};
\node[block, left=of data] (train) {Entra\^inement\\(DVC repro)};
\node[decision, left=of train] (gate) {Seuils\\OK~?};
\node[block, below=of gate] (staging) {D\'eploiement\\staging};
\node[block, right=of staging] (ab) {A/B Test /\\Shadow};
\node[decision, right=of ab] (validate) {M\'etriques\\OK~?};
\node[block, below=of validate] (prod) {Production};

% Arrows
\draw[arrow] (commit) -- (lint);
\draw[arrow] (lint) -- (test);
\draw[arrow] (test) -- (data);
\draw[arrow] (data) -- (train);
\draw[arrow] (train) -- (gate);
\draw[arrow] (gate) -- node[left, label] {Oui} (staging);
\draw[arrow] (staging) -- (ab);
\draw[arrow] (ab) -- (validate);
\draw[arrow] (validate) -- node[right, label] {Oui} (prod);

% Rejection paths
\draw[arrow, dashed, red!60] (gate.north) -- ++(0,0.8) node[above, font=\scriptsize, text=red!70] {Non -- bloquer};
\draw[arrow, dashed, red!60] (validate.east) -- ++(1.2,0) node[right, font=\scriptsize, text=red!70] {Non -- rollback};

\end{tikzpicture}
\caption{Pipeline CI/CD complet pour un projet de Machine Learning.}
\label{fig:cicd-pipeline}
\end{figure}

%% ----------------------------------------------------------------
\section{Comparaison des outils CI/CD}
\label{sec:cicd-comparison}

\begin{formules}[title=\textbf{GitHub Actions vs GitLab CI vs Jenkins vs CircleCI}]
\begin{center}
\begin{tabular}{lcccc}
\toprule
\textbf{Crit\`ere} & \textbf{GitHub Actions} & \textbf{GitLab CI} & \textbf{Jenkins} & \textbf{CircleCI} \\
\midrule
H\'ebergement & Cloud & Cloud/Self & Self-hosted & Cloud \\
Configuration & YAML & YAML & Groovy/YAML & YAML \\
GPU natif & Non\textsuperscript{*} & Oui & Oui & Non\textsuperscript{*} \\
Cache int\'egr\'e & Oui & Oui & Plugin & Oui \\
Marketplace & Tr\`es riche & Mod\'er\'e & Tr\`es riche & Mod\'er\'e \\
Gratuit (public) & Oui & Oui & Oui & Limit\'e \\
Complexit\'e & Faible & Faible & \'Elev\'ee & Faible \\
Int\'egration ML & Bonne & Bonne & Manuelle & Mod\'er\'ee \\
\bottomrule
\end{tabular}
\end{center}
\textsuperscript{*} Possible via des runners auto-h\'eberg\'es.
\end{formules}

\begin{antipattern}[title=\textbf{Erreurs fr\'equentes en CI/CD ML}]
\begin{itemize}
    \item \textbf{Stocker les donn\'ees dans Git}~: les fichiers volumineux ralentissent le d\'ep\^ot et augmentent les co\^uts. Utilisez DVC ou Git LFS.
    \item \textbf{Entra\^iner sur la CI sans cache}~: chaque pipeline r\'e-entra\^ine depuis z\'ero. Utilisez le cache DVC.
    \item \textbf{Ignorer les tests de donn\'ees}~: tester uniquement le code ne suffit pas en ML.
    \item \textbf{D\'eployer sans \emph{gate}}~: l'absence de seuils de performance peut mettre un mod\`ele d\'egrad\'e en production.
    \item \textbf{Secrets en dur}~: ne jamais mettre les cl\'es API dans le code. Utilisez les \emph{secrets} du fournisseur CI.
\end{itemize}
\end{antipattern}

%% ----------------------------------------------------------------
\section{Mini-projet~: mettre en place le CI/CD}
\label{sec:cicd-miniproject}

\begin{miniprojet}[title=\textbf{CI/CD pour le projet de cours}]
\textbf{Objectif}~: configurer un pipeline CI/CD complet pour le projet ML du cours.

\textbf{\'Etapes}~:
\begin{enumerate}
    \item \textbf{Pre-commit}~: cr\'eer le fichier \file{.pre-commit-config.yaml} avec Black, Ruff et mypy. Ex\'ecuter \cli{pre-commit run --all-files} et corriger les erreurs.
    \item \textbf{Tests}~: \'ecrire au moins 5 tests unitaires (fonctions de pr\'etraitement, chargement des donn\'ees, forme des sorties du mod\`ele).
    \item \textbf{Validation}~: cr\'eer le script \file{scripts/validate\_model.py} avec des seuils de performance.
    \item \textbf{GitHub Actions}~: cr\'eer \file{.github/workflows/ml-pipeline.yml} avec les jobs~: lint, test, train, validate.
    \item \textbf{DVC}~: int\'egrer \cli{dvc repro} dans le pipeline CI.
    \item \textbf{Badge}~: ajouter le badge de statut CI dans le \file{README.md}.
\end{enumerate}

\textbf{Livrables}~:
\begin{itemize}
    \item D\'ep\^ot GitHub avec pipeline CI fonctionnel (badge vert).
    \item Capture d'\'ecran de l'ex\'ecution r\'eussie du workflow.
    \item Document expliquant chaque \'etape du pipeline.
\end{itemize}
\end{miniprojet}

%% ----------------------------------------------------------------
\section{Exercices}
\label{sec:cicd-exercises}

\begin{exercice}[Niveau 1 -- Premiers pas avec GitHub Actions]
Cr\'eez un workflow GitHub Actions minimal qui~:
\begin{enumerate}
    \item Se d\'eclenche sur chaque push vers la branche \texttt{main}.
    \item Installe Python 3.11 et les d\'ependances du projet.
    \item Ex\'ecute \cli{ruff check src/} et \cli{pytest tests/}.
\end{enumerate}
V\'erifiez que le workflow s'ex\'ecute correctement en poussant un commit.
\end{exercice}

\begin{exercice}[Niveau 2 -- Validation automatis\'ee du mod\`ele]
\'Etendez le pipeline CI pour inclure une \'etape de validation du mod\`ele~:
\begin{enumerate}
    \item Apr\`es l'entra\^inement, ex\'ecutez un script qui v\'erifie que l'accuracy d\'epasse 0.80 et que le F1-score d\'epasse 0.75.
    \item Ajoutez une v\'erification de drift via le test de Kolmogorov-Smirnov entre les donn\'ees de r\'ef\'erence et les donn\'ees courantes.
    \item Si un seuil n'est pas atteint, le pipeline doit \'echouer avec un message explicite.
    \item Utilisez \cli{dvc metrics diff} pour comparer les m\'etriques avec la branche principale et afficher la diff\'erence dans un commentaire de PR.
\end{enumerate}
\end{exercice}

\begin{exercice}[Niveau 3 -- Pipeline CI/CD de bout en bout avec A/B testing]
Concevez et impl\'ementez un pipeline CI/CD complet~:
\begin{enumerate}
    \item \textbf{CI}~: lint, tests, entra\^inement, validation (seuils + drift + latence).
    \item \textbf{CD staging}~: d\'eploiement automatique sur un environnement de staging via Docker + Kubernetes.
    \item \textbf{A/B testing}~: configurez Istio ou un proxy pour router 10\,\% du trafic vers le nouveau mod\`ele pendant 24h.
    \item \textbf{Promotion}~: si les m\'etriques m\'etier (taux de conversion, latence p95) sont meilleures, promouvoir automatiquement le mod\`ele en production.
    \item \textbf{Rollback}~: impl\'ementez un m\'ecanisme de rollback automatique si les m\'etriques se d\'egradent.
\end{enumerate}
Documentez l'architecture compl\`ete avec un diagramme et justifiez vos choix techniques.
\end{exercice}

%% ----------------------------------------------------------------
\section{Aide-m\'emoire}
\label{sec:cicd-cheatsheet}

\begin{cheatsheet}[title=\textbf{CI/CD pour le Machine Learning}]
\textbf{Concepts cl\'es}~:
\begin{itemize}
    \item \textbf{CI} = int\'egration continue~: tests automatiques \`a chaque commit
    \item \textbf{CD} = livraison/d\'eploiement continu~: mise en production automatis\'ee
    \item \textbf{Gate} = seuil de qualit\'e \`a franchir avant d\'eploiement
\end{itemize}

\textbf{Commandes essentielles}~:
\begin{minted}{bash}
# Pre-commit
pre-commit install          # Activer les hooks
pre-commit run --all-files  # Executer sur tout

# DVC dans la CI
dvc pull                    # Telecharger les donnees
dvc repro                   # Reproduire le pipeline
dvc metrics diff main       # Comparer les metriques

# GitHub Actions (debug local)
act -j test                 # Executer un job localement
\end{minted}

\textbf{Checklist CI/CD ML}~:
\begin{enumerate}
    \item[$\square$] Pre-commit hooks configur\'es (Black, Ruff, mypy)
    \item[$\square$] Tests unitaires avec couverture $\geq 80\,\%$
    \item[$\square$] Validation des donn\'ees (sch\'ema, drift)
    \item[$\square$] Seuils de performance d\'efinis et v\'erifi\'es
    \item[$\square$] V\'erification de la latence d'inf\'erence
    \item[$\square$] Pipeline DVC reproductible
    \item[$\square$] Secrets stock\'es de mani\`ere s\'ecuris\'ee
    \item[$\square$] Strat\'egie de rollback d\'efinie
\end{enumerate}
\end{cheatsheet}
